% from ega3-1-fr.tex
\oldpage[0\textsubscript{III}]{23}
résultent alors de l'hypothèse de récurrence et du lemme
(\hyperref[0.10.3.1.3-fr]{10.3.1.3}). Cette construction faite, il est clair
que l'anneau $B=A_\gamma$ vérifie l'énoncé de
(\hyperref[0.10.3.1-fr]{10.3.1}).
\end{env}

On notera qu'en vertu de (\hyperref[0.10.2.1-fr]{10.2.1, $c)$}), on a un
isomorphisme canonique
\[
\label{0.10.3.1.5-fr}
  \operatorname{gr}(A)\otimes_k K\xrightarrow{\sim}\operatorname{gr}(B).
  \tag{10.3.1.5}
\]

D'autre part, on peut remplacer $B$ par son complété $\mathfrak J B$-adique
$\widehat B$ sans changer les conclusions de
(\hyperref[0.10.3.1-fr]{10.3.1}), puisque $\widehat B$ est un $B$-module plat
(\hyperref[0.7.3.3-fr]{$0_{\mathrm I}$, 7.3.3}), donc un $A$-module plat
(\hyperref[0.6.2.1-fr]{$0_{\mathrm I}$, 6.2.1}).

On a en outre démontré le

\begin{corollary}[10.3.2]
\label{0.10.3.2-fr}
Si $K$ est une extension de degré fini, on peut supposer que $B$ est une
$A$-algèbre finie.
\end{corollary}

\section{Compléments d'algèbre homologique}
\label{section:0.11-fr}

\subsection{Rappels sur les suites spectrales}
\label{subsection:0.11.1-fr}

\begin{env}[11.1.1]
\label{0.11.1.1-fr}
Nous utiliserons dans la suite une notion de suite spectrale plus générale que
celle définie dans (T, 2.4)~; gardant les notations de (T, 2.4), nous
appellerons \emph{suite spectrale} dans une catégorie abélienne $\mathcal C$ un
système E formé des éléments suivants~:
\begin{enumerate}
  \item[(a)] Une famille $(E_r^{pq})$ d'objets de $\mathcal C$ définis pour
    $p\in\mathbb Z$, $q\in\mathbb Z$ et $r\geq2$.
  \item[(b)] Une famille de morphismes
    $d_r^{pq}:E_r^{pq}\to E_r^{p+r,q-r+1}$ tels que
    $d_r^{p+r,q-r+1}d_r^{pq}=0$. On pose
    $Z_{r+1}(E_r^{pq})=\operatorname{Ker}(d_r^{pq})$,
    $B_{r+1}(E_r^{pq})=\operatorname{Im}(d_r^{p-r,q+r-1})$, de sorte que
    \[
      B_{r+1}(E_r^{pq})\subset Z_{r+1}(E_r^{pq})\subset E_r^{pq}.
    \]
  \item[(c)] Une famille d'isomorphismes
    $\alpha_r^{pq}:Z_{r+1}(E_r^{pq})/B_{r+1}(E_r^{pq})
    \xrightarrow{\sim}E_{r+1}^{pq}$.
\end{enumerate}

On définit alors pour $k\geq r+1$, par récurrence sur $k$, les sous-objets
$B_k(E_r^{pq})$ et $Z_k(E_r^{pq})$ comme images réciproques, par le morphisme
canonique $E_r^{pq}\to E_r^{pq}/B_{r+1}(E_r^{pq})$, des sous-objets de ce
quotient identifié par $\alpha_r^{pq}$ aux sous-objets $B_k(E_{r+1}^{pq})$ et
$Z_k(E_{r+1}^{pq})$ respectivement. Il est clair que l'on a alors, à un
isomorphisme près
\[
\label{0.11.1.1.1-fr}
  Z_k(E_r^{pq})/B_k(E_r^{pq})=E_k^{pq}
  \qquad\text{pour }k\geq r+1,
  \tag{11.1.1.1}
\]
et, si on pose encore $B_r(E_r^{pq})=0$ et $Z_r(E_r^{pq})=E_r^{pq}$, on a
les relations d'inclusion
\[
\label{0.11.1.1.2-fr}
  0=B_r(E_r^{pq})\subset B_{r+1}(E_r^{pq})\subset B_{r+2}(E_r^{pq})
  \subset\cdots\subset Z_{r+2}(E_r^{pq})\subset Z_{r+1}(E_r^{pq})
  \subset Z_r(E_r^{pq})=E_r^{pq}.
  \tag{11.1.1.2}
\]

Les autres éléments de la donnée de E sont alors~:
\begin{enumerate}
  \item[(d)] Deux sous-objets $B_\infty(E_2^{pq})$ et
    $Z_\infty(E_2^{pq})$ de $E_2^{pq}$ tels que l'on ait
    $B_\infty(E_2^{pq})\subset Z_\infty(E_2^{pq})$ et, pour tout $k\geq2$,
    \[
      B_k(E_2^{pq})\subset B_\infty(E_2^{pq})
      \qquad\text{et}\qquad
      Z_\infty(E_2^{pq})\subset Z_k(E_2^{pq}).
    \]
    On pose
    \[
    \label{0.11.1.1.3-fr}
      E_\infty^{pq}=Z_\infty(E_2^{pq})/B_\infty(E_2^{pq}).
      \tag{11.1.1.3}
    \]
  \item[(e)]
    